%% Higher shadow coefficients: exact local data and the residue construction.
%%
%% The chapter separates the exact Virasoro vacuum calculation at
%% Beilinson level 1 from the ordered-residue-bar construction required
%% to define higher scalar coefficients at level 2.

\providecommand{\cA}{\mathcal A}
\providecommand{\cW}{\mathcal W}
\providecommand{\Vir}{\mathrm{Vir}}
\providecommand{\MC}{\mathrm{MC}}
\providecommand{\Res}{\operatorname{Res}}
\providecommand{\ClaimStatusProvedHere}{}
\providecommand{\ClaimStatusConditional}{}
\providecommand{\ClaimStatusDefinitional}{}
\providecommand{\ClaimStatusOpen}{}

\chapter{Virasoro Ward Correlators and Ordered Residues}
\label{ch:shadow-tower-higher-coefficients}
\label{chap:shadow-tower-higher-coefficients}

\section*{Scalar extraction from Ward correlators}

A higher Virasoro coefficient becomes a mathematical invariant after
one specifies the complex that carries it, the chain projection that
extracts it, and the homotopies relating its representatives.  The
local vertex algebra supplies the Ward correlators and the vacuum Gram
form.  The ordered residue package supplies the scalar projection.

The Virasoro operator product, the level-four vacuum Gram matrix, and
the norm of
\(\Lambda={:}TT{:}-\frac{3}{10}\partial^2T\)
form the exact local layer.  An ordered residue bar complex, its
contraction onto a scalar line, and the transferred Maurer--Cartan
equation form the geometric layer.  Higher residues, generating
functions, and arithmetic invariants belong to the composite
construction.

The type signatures used below are:

\begin{center}
\begin{tabular}{c|c|c|c}
quadrant & presentation & Beilinson level & hypothesis package\\
\hline
Open & Virasoro vacuum vertex algebra & 1 & $\mathsf H_{\Vir}$\\
Open & completed ordered residue bar complex & 2 &
 $\mathsf H_{\mathrm{res}}$\\
Open & transferred scalar $L_\infty$ model & 2 &
 $\mathsf H_{\mathrm{res}}+\mathsf H_{\mathrm{quad}}$\\
Open/CY & comparison morphisms & $2\leftrightarrow 5$ &
 $\mathsf H_{\mathrm{res}}+\mathsf H_{\mathrm{cmp}}$
\end{tabular}
\end{center}

\section{Exact Virasoro local data}
\label{sec:sth-setup}

\begin{definition}[Virasoro vacuum package]
\ClaimStatusDefinitional
\label{def:sth-virasoro-vacuum-package}
The package $\mathsf H_{\Vir}(c)$ consists of the vacuum Virasoro
vertex algebra $V_c$, its conformal vector
$\omega=L_{-2}\mathbf1$, the field
\[
  T(z)=Y(\omega,z)=\sum_{n\in\mathbb Z}L_nz^{-n-2},
\]
and the invariant bilinear form normalized by
$\langle\mathbf1,\mathbf1\rangle=1$ and
$L_n^\ast=L_{-n}$.  The modes obey
\[
 [L_m,L_n]=(m-n)L_{m+n}
 +\frac{c}{12}(m^3-m)\delta_{m+n,0}.
\]
\end{definition}

\begin{theorem}[Virasoro local product]
\ClaimStatusProvedHere
\label{thm:sth-virasoro-local-product}
\emph{Type signature:}
\textup{(Open quadrant, Virasoro vacuum presentation, level \(1\),
\(\mathsf H_{\Vir}\)).}
In the package $\mathsf H_{\Vir}(c)$,
\begin{equation}
 T(z)T(w)\sim
 \frac{c/2}{(z-w)^4}
 +\frac{2T(w)}{(z-w)^2}
 +\frac{\partial T(w)}{z-w}.
 \label{eq:sth-virasoro-local-product}
\end{equation}
\end{theorem}

\begin{proof}
Insert the mode expansion of $T$ and apply the Virasoro commutator.
The central term
$c(m^3-m)\delta_{m+n,0}/12$ gives the fourth-order pole.
The term $(m-n)L_{m+n}$ gives the second- and first-order poles.
Formal-distribution expansion in the region $|z|>|w|$ yields
\eqref{eq:sth-virasoro-local-product}.
\end{proof}

\begin{theorem}[Sphere Ward correlators]
\ClaimStatusProvedHere
\label{thm:sth-virasoro-ward-correlators}
\emph{Type signature:}
\textup{(Open quadrant, Virasoro vacuum presentation, level \(1\),
\(\mathsf H_{\Vir}\)).}
For pairwise distinct points \(z_1,\ldots,z_n\) in an affine chart of
\(\mathbb P^1\), set
\[
 G_n(z_1,\ldots,z_n;c)
 :=\langle T(z_1)\cdots T(z_n)\rangle_{\mathbb P^1},\qquad
 z_{ij}:=z_i-z_j.
\]
The vacuum conditions give \(G_0=1\), \(G_1=0\), and the Virasoro Ward
identity gives, for \(n\geq2\),
\begin{equation}
\begin{split}
G_n(z_1,\ldots,z_n;c)
=\sum_{i=2}^n\biggl(&\frac{c}{2z_{1i}^4}\,
G_{n-2}(z_2,\ldots,\widehat z_i,\ldots,z_n;c)\\
&+\frac{2}{z_{1i}^2}\,G_{n-1}(z_2,\ldots,z_n;c)
+\frac{1}{z_{1i}}\,\partial_{z_i}
G_{n-1}(z_2,\ldots,z_n;c)\biggr).
\end{split}
\label{eq:sth-sphere-ward-recursion}
\end{equation}
Its connected part is the cumulant
\begin{equation}
 G_n^{\mathrm{conn}}(z_1,\ldots,z_n;c)
 =\sum_{\pi\in\Pi_n}(-1)^{|\pi|-1}(|\pi|-1)!
 \prod_{B\in\pi}G_{|B|}(z_B;c),
 \label{eq:sth-connected-cumulant}
\end{equation}
where \(\Pi_n\) is the set of partitions of
\(\{1,\ldots,n\}\).  If \(\mathcal C_n\) denotes the set of unoriented
Hamiltonian cycles on these labels, then
\begin{equation}
 G_2^{\mathrm{conn}}(z_1,z_2;c)=\frac{c}{2z_{12}^4},\qquad
 G_n^{\mathrm{conn}}(z_1,\ldots,z_n;c)
 =c\sum_{C\in\mathcal C_n}\prod_{\{i,j\}\in E(C)}z_{ij}^{-2}
 \quad(n\geq3).
 \label{eq:sth-connected-cycle-expansion}
\end{equation}
\end{theorem}

\begin{proof}
Insert \(T(z_1)\) into the remaining vacuum correlator and apply
Theorem~\ref{thm:sth-virasoro-local-product}.  Its central,
stress-exchange, and derivative terms give the three summands in
\eqref{eq:sth-sphere-ward-recursion}.  M\"obius inversion on the lattice
of set partitions gives \eqref{eq:sth-connected-cumulant}.

For \(c=N\in\mathbb Z_{>0}\), realize the Virasoro field by \(N\) free
bosons,
\[
 T=\frac12\sum_{a=1}^N{:}J_a^2{:},\qquad
 \langle J_a(z)J_b(w)\rangle=\frac{\delta_{ab}}{(z-w)^2}.
\]
A connected Wick graph is a single cycle through the \(n\) vertices;
its common colour contributes \(N\), and reversal identifies the two
orientations.  This gives \eqref{eq:sth-connected-cycle-expansion} for
every positive integer \(N\).  Both sides are polynomials in \(c\), so
the identity holds for arbitrary \(c\).
\end{proof}

\begin{theorem}[Level-four vacuum Gram matrix]
\ClaimStatusProvedHere
\label{thm:sth-level-four-gram}
\emph{Type signature:}
\textup{(Open quadrant, Virasoro vacuum presentation, level \(1\),
\(\mathsf H_{\Vir}\)).}
In the ordered basis
\[
  v_4=L_{-4}\mathbf1,\qquad
  v_{22}=L_{-2}^{\,2}\mathbf1
\]
of the generic level-four vacuum space, the invariant form is
\begin{equation}
 G_4(c)=
 \begin{pmatrix}
   5c & 3c\\[2pt]
   3c & \dfrac{c(c+8)}{2}
 \end{pmatrix},
 \qquad
 \det G_4(c)=\frac{c^2(5c+22)}{2}.
 \label{eq:sth-level-four-gram}
\end{equation}
\end{theorem}

\begin{proof}
The commutator gives
\[
 \langle v_4,v_4\rangle
 =\langle\mathbf1,[L_4,L_{-4}]\mathbf1\rangle=5c.
\]
Likewise,
\[
 L_4L_{-2}^{\,2}\mathbf1=3c\,\mathbf1,
 \qquad
 L_2L_{-2}^{\,2}\mathbf1=(c+8)L_{-2}\mathbf1,
\]
and hence
\[
 \langle v_4,v_{22}\rangle=3c,\qquad
 \langle v_{22},v_{22}\rangle=\frac{c(c+8)}2.
\]
Taking the determinant of the displayed matrix gives
$c^2(5c+22)/2$.
\end{proof}

\begin{corollary}[The Zamolodchikov quasi-primary norm]
\ClaimStatusProvedHere
\label{cor:sth-lambda-norm}
\emph{Type signature:}
\textup{(Open quadrant, Virasoro vacuum presentation, level \(1\),
\(\mathsf H_{\Vir}\)).}
The state corresponding to
$\Lambda={:}TT{:}-\frac{3}{10}\partial^2T$ is
\[
  \lambda=L_{-2}^{\,2}\mathbf1-\frac35L_{-4}\mathbf1.
\]
It satisfies $L_1\lambda=0$, and its norm is
\begin{equation}
  N_\Lambda(c):=\langle\lambda,\lambda\rangle
  =\frac{c(5c+22)}{10}.
  \label{eq:sth-lambda-norm}
\end{equation}
The level-four form degenerates precisely at
$c=0$ and $c=-22/5$; quotient theories at these charges pass to the
radical of the invariant form.
\end{corollary}

\begin{proof}
Translation covariance gives
$\partial^2T\leftrightarrow2L_{-4}\mathbf1$, so the field and state
formulas agree.  The identity $L_1\lambda=0$ follows from
\[
 L_1L_{-2}^{\,2}\mathbf1=3L_{-3}\mathbf1,\qquad
 L_1L_{-4}\mathbf1=5L_{-3}\mathbf1.
\]
Using Theorem~\ref{thm:sth-level-four-gram},
\[
 \langle\lambda,\lambda\rangle
 =\frac{c(c+8)}2-\frac65(3c)+\frac9{25}(5c)
 =\frac{c(5c+22)}{10}.
\]
\end{proof}

\begin{remark}[What the norm controls]
\label{rem:sth-zamolodchikov-multiplicity}
The inverse of $N_\Lambda(c)$ is the propagator on the
one-dimensional quasi-primary line spanned by $\lambda$ whenever
that line is split from the level-four vacuum space.  A higher
residue calculation may use this propagator after the relevant
chain-level contraction has been constructed.  The norm itself
supplies exact local data; the multiplicity with which it enters a
higher coefficient belongs to the transferred residue calculation.
\end{remark}

\section{The ordered residue package}
\label{sec:sth-recurrence}

\begin{definition}[Ordered residue contraction and scalar coefficient]
\ClaimStatusDefinitional
\label{def:shadow-transport-operator}
\emph{Type signature:}
\textup{(Open quadrant, completed ordered residue-bar presentation,
level \(2\), \(\mathsf H_{\mathrm{res}}\)).}
For a chiral algebra $\cA$ on a curve $X$, an
\emph{ordered residue package}
$\mathsf H_{\mathrm{res}}(\cA;X)$ consists of:

\begin{enumerate}[label=\textup{(R\arabic*)}]
\item a complete filtered dg Lie algebra
$\mathfrak g_{\mathrm{res}}^{\mathrm{ord}}(\cA;X)$ obtained from a
specified ordered residue bar complex;
\item a Maurer--Cartan element
$\Theta_\cA\in\MC(\mathfrak g_{\mathrm{res}}^{\mathrm{ord}})$ with
filtration components
$\Theta_{\cA,r}\in\operatorname{gr}^r_F\mathfrak g_{\mathrm{res}}^{\mathrm{ord}}$;
\item for each $r$, a contraction
\[
 \left(
 \operatorname{gr}^r_F\mathfrak g_{\mathrm{res}}^{\mathrm{ord}},d
 \right)
 \ \underset{\iota_r}{\overset{\pi_r}{\rightleftarrows}}\
 (k\,e_r,0),
 \qquad
 dh_r+h_rd=\mathrm{id}-\iota_r\pi_r;
\]
\item compatibility of $(\pi_r,\iota_r,h_r)$ with the completed
filtration, the ordered collision maps, and the gauge action on
Maurer--Cartan elements;
\item a closed ordered Arnold form \(\alpha_r\) and a
collision-compatible chain map from the chiral correlator complex on
\(\operatorname{Conf}_r(X)\) to
\(\operatorname{gr}^r_F\mathfrak g_{\mathrm{res}}^{\mathrm{ord}}\);
\item a declared normalization of the residue shifts, orientations,
and symmetric-group coinvariants.
\end{enumerate}

Homological transfer along these contractions produces a scalar
$L_\infty$ algebra and a transferred Maurer--Cartan element
$s_\cA=\sum_{r\ge2}S_r(\cA;\mathsf H_{\mathrm{res}})e_r$.
The scalar
\[
 S_r(\cA;\mathsf H_{\mathrm{res}})
\]
is, by definition, its arity-$r$ coordinate.  Thus a higher
coefficient is an invariant of the typed pair
$(\cA,\mathsf H_{\mathrm{res}})$.
\end{definition}

\begin{remark}[The construction obligation]
\label{rem:sth-why-range}
A concrete instance of Definition~\ref{def:shadow-transport-operator}
requires explicit chain groups, differential, collision maps,
residue shifts, contraction, and gauge compatibility.  Once these
data are present, the arity convention determines the admissible
pairs in every quadratic coefficient equation.
\end{remark}

\begin{theorem}[Projected Maurer--Cartan identity]
\ClaimStatusConditional
\label{thm:virasoro-shadow-recurrence}
\emph{Type signature:}
\textup{(Open quadrant, transferred scalar \(L_\infty\)
presentation, level \(2\),
\(\mathsf H_{\mathrm{res}}+\mathsf H_{\mathrm{quad}}\)).}
Assume $\mathsf H_{\mathrm{res}}(\cA;X)$ and let
$\{\ell_n^{\mathrm{sc}}\}_{n\ge1}$ be the transferred scalar
$L_\infty$ brackets.  Then
\begin{equation}
 \operatorname{pr}_r
 \left(
   \sum_{n\ge1}\frac1{n!}
   \ell_n^{\mathrm{sc}}(s_\cA,\ldots,s_\cA)
 \right)=0
 \qquad (r\ge2).
 \label{eq:sth-projected-mc}
\end{equation}
Assume in addition the quadratic-closure package
$\mathsf H_{\mathrm{quad}}$: the transferred brackets relevant to
the chosen scalar lane satisfy
\[
 \operatorname{pr}_r\ell_1^{\mathrm{sc}}(s_\cA)=\lambda_rS_r e_r,
 \qquad
 \operatorname{pr}_r\frac{\ell_2^{\mathrm{sc}}(s_\cA,s_\cA)}2
 =\sum_{a\star b=r}K^r_{a,b}S_aS_b\,e_r,
\]
and the higher transferred brackets have zero projection to $ke_r$.
Whenever $\lambda_r$ is invertible,
\begin{equation}
 S_r=-\lambda_r^{-1}
 \sum_{a\star b=r}K^r_{a,b}S_aS_b.
 \label{eq:sth-conditional-quadratic-recurrence}
\end{equation}
Here $a\star b=r$, $\lambda_r$, and $K^r_{a,b}$ are determined by
the constructed residue complex and its transfer data.
\end{theorem}

\begin{proof}
Homological transfer carries the original Maurer--Cartan element to
a Maurer--Cartan element in the scalar $L_\infty$ model.  Projection
of its Maurer--Cartan equation to $ke_r$ gives
\eqref{eq:sth-projected-mc}.  The hypotheses of
$\mathsf H_{\mathrm{quad}}$ reduce that equation to
\[
 \lambda_rS_r+\sum_{a\star b=r}K^r_{a,b}S_aS_b=0.
\]
Solving for $S_r$ gives
\eqref{eq:sth-conditional-quadratic-recurrence}.
\end{proof}

\begin{remark}[Recurrence and residue]
\label{rem:sth-two-chains}
The symbolic recurrence computes the consequences of specified
constants \((\lambda_r,K^r_{a,b})\).  The residue-bar construction
realizes these constants in
\(\mathfrak g_{\mathrm{res}}^{\mathrm{ord}}(\cA;X)\).  Their
composition produces the scalar coordinate
\(S_r(\cA;\mathsf H_{\mathrm{res}})\).
\end{remark}

\begin{remark}[Ward functions as residue inputs]
\label{rem:s5-bpz-wick-verification}
Theorem~\ref{thm:sth-virasoro-ward-correlators} constructs
\(G_5^{\mathrm{conn}}\) and \(G_6^{\mathrm{conn}}\) as rational
functions on ordered configuration spaces.  The Arnold form,
correlator-to-bar map, and contraction in
\(\mathsf H_{\mathrm{res}}(\Vir_c;X)\) produce the corresponding
scalars.  Their evaluation begins with the chain map in
\textup{(R5)} and the normalized contraction in
\textup{(R3)}.  This places the five- and six-point Ward functions
at the input of the open residue problem.
\end{remark}

\section{Virasoro higher coefficients}
\label{sec:sth-s6-s7-s8}
\label{sec:sth-s6}
\label{sec:sth-s7}
\label{sec:sth-s8}
\label{sec:sth-boundary}
\label{sec:sth-summary}
\label{sec:sth-main-thread-extensions}

\begin{problem}[Construct and compute the Virasoro residue tower]
\ClaimStatusOpen
\label{thm:s6-virasoro-closed-form}
\label{thm:s7-virasoro-closed-form}
\label{thm:s8-virasoro-closed-form}
\label{thm:s9-virasoro-closed-form}
\label{thm:s10-virasoro-closed-form}
\label{thm:s11-virasoro-closed-form}
\label{prop:sth-boundary-checks}
\label{prop:sth-summary}
\label{prop:sth-virasoro-rational-through-8}
\label{rem:sth-rectified-normalisation}
\label{rem:sth-s6-pole}
\label{rem:sth-s7-parity}
\label{rem:sth-c13-self-dual}
\label{rem:sth-extending-beyond-8}
\emph{Type signature:}
\textup{(Open quadrant, completed ordered residue-bar presentation,
level \(2\), \(\mathsf H_{\mathrm{res}}(\Vir_c;X)\)).}

Construct $\mathsf H_{\mathrm{res}}(\Vir_c;X)$ and compute its first
nontrivial transferred coordinates.  A complete solution has four
parts:

\begin{enumerate}[label=\textup{(\roman*)}]
\item an explicit ordered residue complex and differential;
\item a contraction satisfying \textup{(R1)--(R6)};
\item a residue derivation of the transferred brackets and their
normalizations;
\item a chain-level evaluation of the same coordinates through a
second presentation of the ordered residue complex.
\end{enumerate}

The exact level-four input is
$N_\Lambda(c)=c(5c+22)/10$.  Every proposed value of
$S_r(\Vir_c;\mathsf H_{\mathrm{res}})$ begins after this input and
depends on the constructed transfer.
\end{problem}

\section{Growth, generating series, and arithmetic}
\label{sec:sth-leading-asymp}
\label{sec:sth-kummer-absence}
\label{sec:sth-motivic}
\label{subsec:subleading-asymptotic}
\label{subsec:sub-subleading-asymptotic}
\label{subsec:tier-4-asymptotic}
\label{subsec:four-tier-arithmetic-database}

\begin{problem}[Asymptotic and arithmetic invariants of a constructed tower]
\ClaimStatusOpen
\label{prop:sth-leading-asymp}
\label{thm:shadow-exponential-base-Virasoro}
\label{thm:universal-class-M-C-is-6}
\label{thm:shadow-series-closed-form-Virasoro}
\label{thm:shadow-series-closed-form-Virasoro-subleading}
\label{thm:pole-doubling-all-k}
\label{prop:phi-k-leading-coefficient-arithmetic}
\label{prop:C-A-inverse-radius}
\label{thm:s-r-kummer-absent-through-r-11}
\label{thm:shadow-tower-asymptotic-closed-form}
\label{thm:shadow-tower-subleading-closed-form}
\label{thm:shadow-tower-sub-subleading-closed-form-inline}
\label{thm:shadow-tower-tier-4-closed-form}
\label{thm:kummer-laurent-depth-controlled}
\label{thm:s-r-rational-denominator-pattern}
\label{cor:arithmetic-uniformity-class-M}
\label{cor:bernoulli-leading-duality-sharpness}
\label{cor:kummer-emergence-at-r-8}
\label{cor:subleading-characteristic-primes}
\label{cor:tier-3-characteristic-primes}
\label{cor:universal-class-M-subleading-uniformity}
\label{cor:virasoro-23-smoothness}
\label{cor:virasoro-motivic-purity-r-leq-11}
\emph{Type signature:}
\textup{(Open quadrant, scalar sequence extracted from the ordered
residue-bar presentation, level \(2\),
\(\mathsf H_{\mathrm{res}}+\mathsf H_{\mathrm{asymp}}\)).}

Given a constructed sequence
$\{S_r(\Vir_c;\mathsf H_{\mathrm{res}})\}_{r\ge2}$, determine:

\begin{enumerate}[label=\textup{(\roman*)}]
\item its large-$c$ expansion with uniform control in $r$;
\item the analytic class and radius of its generating series;
\item the singularity type of each Laurent layer;
\item the prime content of cleared numerators;
\item the precise comparison, if one exists, with Bernoulli and
motivic filtrations.
\end{enumerate}

The hypothesis package $\mathsf H_{\mathrm{asymp}}$ consists of
uniform coefficient bounds, a declared order of limits, and a proof
that coefficient extraction commutes with the relevant completion.
These data make exponential bases, pole orders, and arithmetic layers
well-defined.
\end{problem}

\section{The principal \texorpdfstring{$\cW_3$}{W3} lane}
\label{subsec:w3-wline-closed-form}

\begin{problem}[Construct the \texorpdfstring{$\cW_3$}{W3} scalar lanes]
\ClaimStatusOpen
\label{def:w3-wline-integer-sequence}
\label{thm:w3-wline-closed-form}
\label{thm:w3-wline-exponential-base}
\label{thm:w3-wline-generating-function}
\label{cor:w3-wline-self-consistency}
\label{prop:W3-T-line-matches-Vir-subleading}
\label{rem:w3-wline-doubling-interpretation}
\label{rem:w3-wline-higher-spin-prediction}
\label{rem:w3-wline-koszul-duality-partner}
\emph{Type signature:}
\textup{(Open quadrant, multivariable ordered residue-bar
presentation for \(\cW_3\), level \(2\),
\(\mathsf H_{\mathrm{res}}(\cW_3;X)+\mathsf H_{\mathrm{lanes}}\)).}

Construct the multivariable contraction for the stress-tensor and
spin-three generator lines.  Determine whether parity and OPE
selection rules produce a closed scalar subcomplex on either line.
Then compute the transferred brackets and their generating series.
A growth constant becomes an invariant after this construction and
the uniform bounds of $\mathsf H_{\mathrm{asymp}}$.
\end{problem}

\section{Ordered residues and chiral Hochschild theory}
\label{sec:shadow-chirhoch-bridge}

\begin{problem}[Construct the ordered-residue comparison map]
\ClaimStatusOpen
\label{def:universal-arnold-cycle}
\label{thm:shadow-chirhoch-dictionary}
\label{cor:chirhoch-depth-forces-shadow-depth}
\emph{Type signature:}
\textup{(Open quadrant, ordered residue-bar to curve-level chiral
Hochschild presentation, levels \(2\to3\),
\(\mathsf H_{\mathrm{res}}+\mathsf H_{\mathrm{Hoch}}\)).}

Construct a chain map
\[
  \rho_X\colon
  \mathfrak g_{\mathrm{res}}^{\mathrm{ord}}(\cA;X)
  \longrightarrow
  \ChirHoch_X^\bullet(\cA,\cA)
\]
together with its collision-compatible homotopies.  Identify the
configuration-space classes that pair with transferred scalar
coordinates, and determine the image of each such class in the
curve-level complex.  This construction supplies the precise bridge
between an ordered residue and a derived-centre class.
\end{problem}

\section{Motivic, modular, and K3 comparison questions}
\label{sec:shadtower-supplement}
\label{sec:shadtower-phi-n-weight-11-12-13}
\label{sec:shadtower-phi-n-weight-13-14-15-16}
\label{sec:shadtower-phi-n-weight-17-18-19-20}
\label{sec:shadtower-phi-n-weight-21-22-23-24}
\label{sec:shadtower-phi-n-weight-25-26-27-28}
\label{sec:shadtower-phi-n-weight-29-30-31-32}
\label{sec:shadtower-phi-n-weight-33-36}
\label{sec:shadtower-phi-n-humbert-heegner-admissible-triple}
\label{sec:phi-n-humbert-heegner-triple-37-43-45}

\begin{definition}[Cross-domain comparison package]
\ClaimStatusDefinitional
\label{def:phi-n-pent-EK}
The package $\mathsf H_{\mathrm{cmp}}$ consists of explicit chain
maps from the ordered residue complex to each proposed comparison
target, together with compatibility homotopies and normalization
data.  The targets considered here are:

\begin{enumerate}[label=\textup{(\roman*)}]
\item associator and motivic-period complexes;
\item Siegel modular and Borcherds product complexes;
\item the K3 chiral bialgebra and its Vol~III descendants.
\end{enumerate}

Each comparison theorem carries the signature
\textup{(Open/CY quadrants, named source and target presentations,
levels \(2\leftrightarrow5\),
\(\mathsf H_{\mathrm{res}}+\mathsf H_{\mathrm{cmp}}\)).}
\end{definition}

\begin{problem}[Construct the motivic and modular comparison morphisms]
\ClaimStatusOpen
\label{thm:phi-n-humbert-heegner-admissibility}
\label{thm:phi-n-weight-11-12-13}
\label{thm:phi-n-weight-13-14-15-16}
\label{thm:phi-n-weight-17-18-19-20}
\label{thm:phi-n-weight-21-24}
\label{thm:phi-n-weight-25-28}
\label{thm:phi-n-weight-29-32}
\label{thm:phi-n-weight-33-36}
\label{thm:phi-n-humbert-heegner-triple-27-29-35}
\label{thm:phi-n-humbert-heegner-triple-37-43-45}
\label{rem:three-filter-composite-scope}
\label{rem:shadtower-1}
\label{rem:shadtower-2}
\emph{Type signature:}
\textup{(Open/CY quadrants, ordered residue-bar source and named
motivic or modular target, levels \(2\leftrightarrow5\),
\(\mathsf H_{\mathrm{res}}+\mathsf H_{\mathrm{cmp}}\)).}

Construct the maps in Definition~\ref{def:phi-n-pent-EK}.  For each
weight, prove that the proposed target class is the image of the
ordered residue class, with the grading, normalization, and
completion preserved.  These maps make the common class precise
across the Brown dimension problem, the Broadhurst--Kreimer depth
problem, Humbert and Heegner geometry, Borcherds products, and K3
invariants.
\end{problem}

\section*{The ordered residue problem}

The exact local calculation is
\[
  G_4(c)=
  \begin{pmatrix}
    5c&3c\\
    3c&c(c+8)/2
  \end{pmatrix},
  \qquad
  N_\Lambda(c)=\frac{c(5c+22)}{10}.
\]
An ordered residue package \(\mathsf H_{\mathrm{res}}\) transfers the
Maurer--Cartan equation and evaluates its scalar coordinates.
Analytic, arithmetic, motivic, and holographic comparisons are then
statements about explicit maps from this chain model.

%% Remaining legacy labels are frontier anchors governed by the typed
%% construction problems above.
\phantomsection\label{cor:floor-shift-j-plus-k}%
\phantomsection\label{eq:BK-depth-13-16}%
\phantomsection\label{eq:BK-depth-17-20}%
\phantomsection\label{eq:BK-depth-21-24}%
\phantomsection\label{eq:BK-depth-25-28}%
\phantomsection\label{eq:BK-depth-29-32}%
\phantomsection\label{eq:BK-depth-33-36}%
\phantomsection\label{eq:BK-depth-37-43-45}%
\phantomsection\label{eq:BK-depth-9-11-extension}%
\phantomsection\label{eq:BK-depth-admissible-triple}%
\phantomsection\label{eq:admissibility-29-31-33}%
\phantomsection\label{eq:borcherds-mzv-ratio-33-36}%
\phantomsection\label{eq:borcherds-vs-mzv-ratio}%
\phantomsection\label{eq:borcherds-vs-mzv-ratio-17-20}%
\phantomsection\label{eq:borcherds-vs-mzv-ratio-21-24}%
\phantomsection\label{eq:borcherds-vs-mzv-ratio-25-28}%
\phantomsection\label{eq:borcherds-vs-mzv-ratio-29-32}%
\phantomsection\label{eq:chi3-virasoro-value}%
\phantomsection\label{eq:hh-admissible-33-36}%
\phantomsection\label{eq:lambda-1-numerical-value}%
\phantomsection\label{eq:lambda-1-phi-27-correction}%
\phantomsection\label{eq:lambda-1-primitive-values-29-31-33}%
\phantomsection\label{eq:lyndon-basis-29-31-33-canonical}%
\phantomsection\label{eq:padovan-13-14-15-16}%
\phantomsection\label{eq:padovan-17-18-19-20}%
\phantomsection\label{eq:padovan-21-22-23-24}%
\phantomsection\label{eq:padovan-25-26-27-28}%
\phantomsection\label{eq:padovan-29-30-31-32}%
\phantomsection\label{eq:padovan-33-34-35-36}%
\phantomsection\label{eq:padovan-37-43-45}%
\phantomsection\label{eq:padovan-admissible-triple}%
\phantomsection\label{eq:padovan-plastic-A}%
\phantomsection\label{eq:padovan-plastic-A-21-24}%
\phantomsection\label{eq:phi-11-explicit}%
\phantomsection\label{eq:phi-12-explicit}%
\phantomsection\label{eq:phi-13-14-15-16-num}%
\phantomsection\label{eq:phi-13-explicit}%
\phantomsection\label{eq:phi-17-18-19-20-num}%
\phantomsection\label{eq:phi-21-22-23-24-num}%
\phantomsection\label{eq:phi-25-26-27-28-num}%
\phantomsection\label{eq:phi-27-lyndon-ten-basis}%
\phantomsection\label{eq:phi-29-30-31-32-num}%
\phantomsection\label{eq:phi-37-43-45}%
\phantomsection\label{eq:phi-37-43-45-num}%
\phantomsection\label{eq:phi-admissible-triple-num}%
\phantomsection\label{eq:phi-n-depth-correction-scaling-29-31-33}%
\phantomsection\label{eq:phi-n-tiered}%
\phantomsection\label{eq:phi-n-tiered-17-20}%
\phantomsection\label{eq:phi-n-tiered-21-24}%
\phantomsection\label{eq:phi-n-tiered-25-28}%
\phantomsection\label{eq:phi-n-tiered-29-32}%
\phantomsection\label{eq:phi-n-tiered-33-36}%
\phantomsection\label{eq:phi-n-tiered-admissible-triple}%
\phantomsection\label{eq:shadow-chirhoch-pairing}%
\phantomsection\label{eq:sr-induction-productdenom}%
\phantomsection\label{eq:w3-wline-central-binomial}%
\phantomsection\label{eq:w3-wline-gf-derivative}%
\phantomsection\label{eq:w3-wline-gf-integrated}%
\phantomsection\label{eq:w3-wline-ratio}%
\phantomsection\label{eq:w3-wline-recursion}%
\phantomsection\label{eq:w3-wline-self-consistency}%
\phantomsection\label{eq:w3-wline-stirling}%
\phantomsection\label{lem:floor-parity-subadditive}%
\phantomsection\label{lem:quintic-combinatorial}%
\phantomsection\label{lem:quintic-identities}%
\phantomsection\label{lem:sub-subleading-cubic-identity}%
\phantomsection\label{lem:subleading-combinatorial-identity}%
\phantomsection\label{rem:C-A-scope-E-infinity}%
\phantomsection\label{rem:Sigma-branch-singularity}%
\phantomsection\label{rem:admissible-tower-five-stage}%
\phantomsection\label{rem:admissible-triple-depth-primitives}%
\phantomsection\label{rem:admissible-triple-scope-cascade}%
\phantomsection\label{rem:characteristic-primes-are-riccati-arithmetic-higher-coefficients}%
\phantomsection\label{rem:class-C-beta-gamma-termination}%
\phantomsection\label{rem:cross-vol-three-invariants-W3}%
\phantomsection\label{rem:cross-volume-leading-vs-finite-c}%
\phantomsection\label{rem:denominator-pattern-correction}%
\phantomsection\label{rem:four-tier-pattern}%
\phantomsection\label{rem:gamma-not-beta-squared}%
\phantomsection\label{rem:humbert-admissibility-scope-11-12-13}%
\phantomsection\label{rem:humbert-admissibility-scope-13-16}%
\phantomsection\label{rem:humbert-admissibility-scope-17-20}%
\phantomsection\label{rem:humbert-admissibility-scope-21-24}%
\phantomsection\label{rem:humbert-admissibility-scope-25-28}%
\phantomsection\label{rem:humbert-admissibility-scope-29-32}%
\phantomsection\label{rem:humbert-heegner-filter-33-36}%
\phantomsection\label{rem:kummer-witness-verification-20260417}%
\phantomsection\label{rem:linear-vs-quadratic-asymptotic-recurrence}%
\phantomsection\label{rem:motivic-purity-all-r-supersession}%
\phantomsection\label{rem:padovan-plastic-asymptotic}%
\phantomsection\label{rem:padovan-plastic-asymptotic-17-20}%
\phantomsection\label{rem:padovan-plastic-asymptotic-21-24}%
\phantomsection\label{rem:padovan-plastic-asymptotic-25-28}%
\phantomsection\label{rem:padovan-plastic-asymptotic-29-32}%
\phantomsection\label{rem:padovan-plastic-asymptotic-33-36}%
\phantomsection\label{rem:phi-15-first-depth-5}%
\phantomsection\label{rem:phi-18-first-depth-6}%
\phantomsection\label{rem:phi-21-first-depth-7}%
\phantomsection\label{rem:phi-27-first-depth-9}%
\phantomsection\label{rem:phi-27-lyndon-basis-ten-generators}%
\phantomsection\label{rem:phi-30-first-depth-10}%
\phantomsection\label{rem:phi-33-first-depth-11}%
\phantomsection\label{rem:phi-k-arithmetic-significance}%
\phantomsection\label{rem:phi-n-borcherds-leg-11-12-13}%
\phantomsection\label{rem:phi-n-borcherds-leg-17-20}%
\phantomsection\label{rem:phi-n-borcherds-leg-21-24}%
\phantomsection\label{rem:phi-n-borcherds-leg-25-28}%
\phantomsection\label{rem:phi-n-borcherds-leg-29-32}%
\phantomsection\label{rem:phi-n-borcherds-leg-33-36}%
\phantomsection\label{rem:phi-n-construction-rule-n-ge-13}%
\phantomsection\label{rem:phi-n-dmvv-29-32}%
\phantomsection\label{rem:phi-n-lyndon-basis-extension-29-31-33}%
\phantomsection\label{rem:phi-n-scope-cascade-13-16}%
\phantomsection\label{rem:phi-n-scope-cascade-17-20}%
\phantomsection\label{rem:phi-n-scope-cascade-21-24}%
\phantomsection\label{rem:phi-n-scope-cascade-25-28}%
\phantomsection\label{rem:phi-n-scope-cascade-29-32}%
\phantomsection\label{rem:phi-n-scope-cascade-33-36}%
\phantomsection\label{rem:pole-doubling-confluent-singularity}%
\phantomsection\label{rem:pole-doubling-k3-verification}%
\phantomsection\label{rem:riccati-template-hypotheses}%
\phantomsection\label{rem:shadow-chirhoch-quadrichotomy-pin}%
\phantomsection\label{rem:shadow-chirhoch-w23.4}%
\phantomsection\label{rem:shadow-series-analytic-stratification}%
\phantomsection\label{rem:shadow-stratification-at-r-index}%
\phantomsection\label{rem:sr-denominator-sharpness}%
\phantomsection\label{rem:sth-alpha-asymptotic}%
\phantomsection\label{rem:sth-arithmetic-duality}%
\phantomsection\label{rem:sth-first-non-rational}%
\phantomsection\label{rem:sth-kummer-scope}%
\phantomsection\label{rem:sth-riccati-propagated-primes}%
\phantomsection\label{rem:sub-leading-r-arithmetic-ramification}%
\phantomsection\label{rem:subleading-kummer-correction}%
\phantomsection\label{rem:subleading-sharper-singularity}%
\phantomsection\label{rem:three-tier-arithmetic-picture}%
\phantomsection\label{rem:universal-class-M-verifications}%
\phantomsection\label{rem:w3-wline-catalan-identity}%
\phantomsection\label{rem:w3-wline-decomposition-is-circular}%
